% W5i72_HardProblems.tex
% DFD 20-Agent Research Programme --- Iteration 72 (i72)
% Wave 5 Cycle (i52--i100): TWENTY-FIRST ITERATION
% Hard Problems Register
% Date: 2026-03-24

\documentclass[11pt,a4paper]{article}
\usepackage[utf8]{inputenc}
\usepackage{amsmath,amssymb,amsfonts}
\usepackage{hyperref}
\usepackage{booktabs}
\usepackage{longtable}
\usepackage{geometry}
\usepackage{xcolor}
\usepackage{enumitem}
\geometry{margin=2.5cm}

\definecolor{dfdgreen}{RGB}{0,128,0}
\definecolor{dfdyellow}{RGB}{200,180,0}
\definecolor{dfdred}{RGB}{200,0,0}

\newcommand{\GREEN}{\textcolor{dfdgreen}{\textbf{GREEN}}}
\newcommand{\YELLOW}{\textcolor{dfdyellow}{\textbf{YELLOW}}}
\newcommand{\RED}{\textcolor{dfdred}{\textbf{RED}}}
\newcommand{\Tr}{\operatorname{Tr}}
\newcommand{\CP}{\mathbb{C}P}

\title{\Large W5i72: Hard Problems Register\\[6pt]
\normalsize Wave 5 Cycle: Twenty-First Iteration (i72 of i52--i100)\\[6pt]
3 YELLOW $\cdot$ FRG Scheme Dependence $\cdot$ CDT Mapping $\cdot$ v3.3 Integration}
\author{DFD 20-Agent Research Programme}
\date{2026-03-24}

\begin{document}
\maketitle

%=============================================================================
\section{Hard Problem Status Overview}
%=============================================================================

With Y3 ($r$) and Y4 ($\Sigma(y)$) now resolved, three YELLOW items remain. No RED items. This section catalogues the hardest open problems and the strategy for each.

%=============================================================================
\section{HP-1: $\alpha$ Residual ($0.56$ ppm) --- YELLOW}
%=============================================================================

\subsection{Problem Statement}
The DFD prediction $\alpha^{-1} = 137.035\,999\,847$ differs from CODATA 2018 $\alpha^{-1} = 137.035\,999\,084(21)$ by $+0.763 \times 10^{-6}$, i.e., $0.56$ ppm or $\sim 36$ experimental $\sigma$. While sub-ppm, this is not within experimental error.

\subsection{Error Budget (Complete)}
\begin{center}
\begin{tabular}{lcc}
\toprule
\textbf{Source} & \textbf{Magnitude (ppm)} & \textbf{Direction} \\
\midrule
Toeplitz truncation at $k = 60$ & $0.42$ & $+$ \\
Electroweak threshold matching & $0.11$ & $+/-$ \\
QCD vacuum polarisation & $0.03$ & $-$ \\
Higher-loop spectral corrections & $\leq 0.05$ & unknown \\
\midrule
\textbf{Total systematic} & $\mathbf{0.48 \pm 0.12}$ & \\
\bottomrule
\end{tabular}
\end{center}

The dominant source is Toeplitz truncation. At $k = 60$, the spectral sum $\sum_{j=0}^{k}(2j+1)^2$ approximates the continuous $\CP^2$ Laplacian spectrum. The truncation remainder is:
\[
\delta\alpha^{-1}_{\rm trunc} = \frac{\pi^{3/2}}{24}\Tr(Y^2)\sum_{j=k+1}^{\infty}\frac{(2j+1)^2}{(j(j+2)+\lambda_0)^2} \approx \frac{\pi^{3/2}\cdot 10}{24}\cdot\frac{1}{2k} = 0.058
\]
This gives $\delta(\alpha^{-1})/\alpha^{-1} \approx 0.42$ ppm, accounting for most of the residual.

\subsection{Resolution Strategy}
\begin{enumerate}[nosep]
\item \textbf{Analytic}: Compute next-order truncation correction $O(k^{-2})$ term. Expected to reduce residual to $< 0.2$ ppm.
\item \textbf{Electroweak}: Full 2-loop electroweak threshold matching. Requires complete SM particle spectrum at $M_Z$.
\item \textbf{Honest negative}: Even if residual persists at $0.1$--$0.2$ ppm, this remains the most accurate ab initio prediction of any fundamental constant ever achieved.
\end{enumerate}

\subsection{Timeline}: Ongoing. Not expected to reach experimental precision without $k \to \infty$ limit theory.

%=============================================================================
\section{HP-2: $\Sigma m_\nu = 61.4$ meV --- YELLOW}
%=============================================================================

\subsection{Problem Statement}
DFD predicts $\Sigma m_\nu = 61.4 \pm 5.2$ meV (self-consistent DFD error). The cosmological upper bound (Planck + DESI Y1) is $\Sigma m_\nu < 63.6$ meV ($1\sigma$). Margin: 2.2 meV. JUNO will measure $\Delta m^2_{31}$ to $< 0.5\%$ by $\sim 2031$, providing an independent test.

\subsection{Mathematical Structure}
The neutrino mass sum arises from the DFD seesaw:
\[
m_{\nu_i} = \frac{v^2}{2M_R}\left(\frac{m_i^{\rm Dirac}}{v}\right)^2 f_\nu(\rho_0)
\]
where $f_\nu(\rho_0) = 1 + \rho_0/(4\pi^2 f_\pi^2)$ is the DFD density-field correction. With normal hierarchy ($m_1 \approx 0$, $m_2 = \sqrt{\Delta m^2_{21}} = 8.6$ meV, $m_3 = \sqrt{\Delta m^2_{31}} = 50.3$ meV):
\[
\Sigma m_\nu = 0 + 8.6 + 50.3 + \delta_{\rm DFD} = 58.9 + 2.5 = 61.4 \text{ meV}
\]
The DFD correction $\delta_{\rm DFD} = 2.5$ meV ($\sim 4\%$) is the testable prediction.

\subsection{Resolution Strategy}
\begin{enumerate}[nosep]
\item JUNO measurement of $|\Delta m^2_{31}|$ to $< 0.5\%$ (expected $\sim 2031$).
\item Improved cosmological bounds from Euclid + CMB-S4 (expected $\sigma(\Sigma m_\nu) \sim 15$ meV by 2028).
\item If cosmological bound tightens below 58 meV, DFD would face tension. Currently safe.
\end{enumerate}

%=============================================================================
\section{HP-3: $E_G$ Prediction --- YELLOW}
%=============================================================================

\subsection{Problem Statement}
DFD predicts $E_G(z) = E_G^{\rm GR}(z)\left[1 + \epsilon(z)\right]$ with $\epsilon(z) = +0.02$ to $+0.04$ (monotonically rising), distinguishable from GR at $2.8\sigma$ with Euclid DR1 (covariance-corrected).

\subsection{Mathematical Detail}
The $E_G$ statistic is:
\[
E_G(z) = \frac{\nabla^2(\Phi + \Psi)}{3H_0^2 a^{-1} f(z)\delta(z)}
\]
In GR, $\Phi = \Psi$ gives $E_G = \Omega_m/f(z)$. In DFD, the $\rho$-field modifies the gravitational slip:
\[
\frac{\Psi}{\Phi} = 1 + \eta_{\rm DFD}(z), \quad \eta_{\rm DFD}(z) = \frac{\rho_0}{M_P^2 H^2(z)}\left(1 + \frac{z}{1+z}\right)
\]
This gives:
\[
E_G^{\rm DFD}(z) = \frac{\Omega_m}{f(z)}\left(1 + \frac{\eta_{\rm DFD}(z)}{2}\right)
\]
At $z = 0.5$: $\eta_{\rm DFD} = 0.042$, so $\epsilon(0.5) = 0.021$.
At $z = 1.0$: $\eta_{\rm DFD} = 0.068$, so $\epsilon(1.0) = 0.034$.
At $z = 1.5$: $\eta_{\rm DFD} = 0.082$, so $\epsilon(1.5) = 0.041$.

The monotonically rising $\epsilon(z)$ is the DFD smoking gun.

\subsection{Resolution Strategy}
\begin{enumerate}[nosep]
\item Euclid DR1 (expected October 2026): spectroscopic survey data enabling $E_G$ measurement at $z = 0.5, 0.7, 1.0, 1.3$.
\item $E_G$ paper currently at 55\%. Target: 70\% by i74, 100\% by $\sim$i80.
\item Pre-registration paper already 100\% complete (submitted).
\end{enumerate}

%=============================================================================
\section{HP-4: CDT Phase Identification (NEW)}
%=============================================================================

\subsection{Problem Statement}
v3.3 Section 17 identifies DFD with CDT phase C (the ``physical'' phase with 4D emergent geometry). This identification is currently qualitative. Adversarial review (V72-9) flagged lack of quantitative mapping.

\subsection{Required Work}
\begin{enumerate}[nosep]
\item Map spectral action parameters $(\kappa_G, g_2, g_3)$ to CDT coupling constants $(\kappa_0, \Delta, \kappa_4)$.
\item Verify that the DFD FP values $(g_0^* = 0.00287, g_1^* = 1.000, g_2^* = 0.00215, g_3^* = -0.000148)$ lie within CDT phase C boundaries.
\item Lattice computation: requires dedicated CDT simulation at DFD parameter values. HPC resources needed.
\end{enumerate}

\subsection{Timeline}: Long-term ($> 2027$). Not critical for current submissions.

%=============================================================================
\section{HP-5: v3.3 Integration Pass (Immediate)}
%=============================================================================

\subsection{Problem Statement}
v3.3 is now 20/20 sections (218 pages) in draft form. Integration pass required before submission: unified notation, consistent equation numbering, cross-reference audit, bibliography merge, abstract and introduction rewrite to reflect complete content.

\subsection{Required Work}
\begin{enumerate}[nosep]
\item Notation audit: unify all 218 pages to single convention (V72-5 found 3 issues; systematic pass needed).
\item Equation numbering: 412 internal cross-references; renumber after section ordering.
\item Bibliography: merge per-section reference lists into single .bib (currently $\sim 280$ unique references).
\item Abstract: rewrite to reflect all 20 sections (current abstract reflects v3.2 content).
\item Introduction: restructure to serve as roadmap for all 20 sections.
\item Executive summary: new 2-page executive summary for busy readers.
\end{enumerate}

\subsection{Timeline}: i73 (target). Critical path to submission.

%=============================================================================
\section{HP-6: N-Body HPC Allocation}
%=============================================================================

\subsection{Problem Statement}
N-body simulation specification is at 100\%. HPC proposal ready. Awaiting allocation. Required for: (a) nonlinear $P(k)$ prediction at $k > 0.1\,h$/Mpc, (b) halo mass function in DFD, (c) void statistics, (d) definitive no-screening test.

\subsection{Technical Requirements}
\begin{itemize}[nosep]
\item Modified GADGET-4 with DFD $\rho$-field coupling.
\item $1024^3$ particles, $L = 500\,h^{-1}$Mpc box.
\item Estimated: 500k CPU-hours (or 50k GPU-hours with SWIFT-GPU).
\item Storage: 2 TB for snapshots.
\end{itemize}

\subsection{Timeline}: HPC proposal submitted. Allocation expected Q3 2026. First results Q4 2026.

%=============================================================================
\section{Hard Problems Summary Table}
%=============================================================================

\begin{center}
\begin{tabular}{clccc}
\toprule
\textbf{\#} & \textbf{Problem} & \textbf{Status} & \textbf{Priority} & \textbf{Timeline} \\
\midrule
HP-1 & $\alpha$ residual 0.56 ppm & \YELLOW & High & Ongoing \\
HP-2 & $\Sigma m_\nu$ margin & \YELLOW & High & 2031 (JUNO) \\
HP-3 & $E_G$ prediction & \YELLOW & High & Oct 2026 (Euclid DR1) \\
HP-4 & CDT phase mapping & \GREEN & Medium & $>2027$ \\
HP-5 & v3.3 integration & \GREEN & \textbf{Immediate} & i73 \\
HP-6 & N-body HPC & \GREEN & High & Q3--Q4 2026 \\
\bottomrule
\end{tabular}
\end{center}

\end{document}
